\ifdefined\gkver\else\def\gkver{student}\fi
\documentclass[\gkver]{gaokaozhenti}
\begin{document}

\chapter{2010年上海}

\begin{enumerate}
\item
%% number: 1
%% paperName: 2010年上海高考第 1 题 2 分
%% typeId: 1
%% type: 单选题
%% chapter: 原子和原子核
%% point: 原子核式结构
%% method: 无
%% score: 2
%% degree: 850
%% duplicateId: 0
%% body:
卢瑟福提出原子的核式结构模型。这一模型建立的基础是\xzanswer{A}
\fourchoices[answer=A]
{$\alpha$ 粒子的散射实验}
{对阴极射线的研究}
{天然放射性现象的发现}
{质子的发现}
%% answer: A
%% memo:
\memoanswer{
卢瑟福提出原子的核式结构模型建立在 $\alpha$ 粒子散射实验的基础上，故 A 正确。
}

\item
%% number: 2
%% paperName: 2010年上海高考第 2 题 2 分
%% typeId: 1
%% type: 单选题
%% chapter: 机械波
%% point: 干涉衍射
%% method: 无
%% score: 2
%% degree: 850
%% duplicateId: 0
%% body:
利用发波水槽得到的水面波形如图~\ref{2010上海02a}、\ref{2010上海02b} 所示，则\xzanswer{D}
\twopicture[labelA=2010上海02a, labelB=2010上海02b, hA=1.7cm, hB=1.7cm]
{figs/02a.png}
{figs/02b.png}
\fourchoices[answer=D]
{图 a、b 均显示了波的干涉现象}
{图 a、b 均显示了波的衍射现象}
{图 a 显示了波的干涉现象，图 b 显示了波的衍射现象}
{图 a 显示了波的衍射现象，图 b 显示了波的干涉现象}
%% answer: D
%% memo:
\memoanswer{
图 a 显示水波遇到障碍物，能够绕到障碍物后面，即应该为衍射图样；图 b 显示两列波相遇后，有些地方振动加强、有些地方振动减弱，即为干涉图样，故 D 正确。
}

\item
%% number: 3
%% paperName: 2010年上海高考第 3 题 2 分
%% typeId: 1
%% type: 单选题
%% chapter: 机械波
%% point: 干涉衍射
%% method: 无
%% score: 2
%% degree: 850
%% duplicateId: 0
%% body:
声波能绕过某一建筑物传播而光波却不能绕过该建筑物，这是因为\xzanswer{C}
\fourchoices[answer=C]
{声波是纵波，光波是横波}
{声波振幅大，光波振幅小}
{声波波长较长，光波波长较短}
{声波波速较小，光波波速很大}
%% answer: C
%% memo:
\memoanswer{
声波的波长比较长，容易发生衍射现象，光波的波长很短，很不容易发生衍射现象，故 C 正确。
}

\item
%% number: 4
%% paperName: 2010年上海高考第 4 题 2 分
%% typeId: 1
%% type: 单选题
%% chapter: 原子和原子核
%% point: 核裂变
%% method: 无
%% score: 2
%% degree: 850
%% duplicateId: 0
%% body:
现已建成的核电站发电的能量来自于\xzanswer{C}
\fourchoices[answer=C]
{天然放射性元素放出的能量}
{人工放射性同位素放出的能量}
{重核裂变放出的能量}
{化学反应放出的能量}
%% answer: C
%% memo:
\memoanswer{
核电站发电的能量主要来自重核裂变放出的能量，故 C 正确。
}

\item
%% number: 5
%% paperName: 2010年上海高考第 5 题 2 分
%% typeId: 1
%% type: 单选题
%% chapter: 电路
%% point: 动态分析
%% method: 无
%% score: 2
%% degree: 650
%% duplicateId: 0
%% body:
在右图的闭合电路中，当滑片 P 向右移动时，两电表读数的变化是\xzanswer{B}
\onepicture[w=4.5cm]{\includesvg{figs/05}}
\fourchoices[answer=B]
{A 变大，V 变大}
{A 变小，V 变大}
{A 变大，V 变小}
{A 变小，V 变小}
%% answer: B
%% memo:
\memoanswer{
当滑片 P 向右移动时，滑动变阻器接入电路的阻值增大，则电路的总电阻增大，由 $I=\frac{E}{R+r+R_{0}}$ 可知，电路的电流 $I$ 减小；由闭合电路欧姆定律 $U=E-I(r+R_{0})$ 可知，电压表的读数增大，故 B 正确。
}

\item
%% number: 6
%% paperName: 2010年上海高考第 6 题 2 分
%% typeId: 1
%% type: 单选题
%% chapter: 光学
%% point: 光子说
%% method: 无
%% score: 2
%% degree: 850
%% duplicateId: 0
%% body:
根据爱因斯坦光子说，光子能量 $E$ 等于（$h$ 为普朗克常量，$c$、$\lambda$ 为真空中的光速和波长）\xzanswer{A}
\fourchoices[answer=A]
{$\frac{hc}{\lambda}$}
{$\frac{h\lambda}{c}$}
{$h\lambda$}
{$\frac{h}{\lambda}$}
%% answer: A
%% memo:
\memoanswer{
由于光子能量为 $E = h\nu$（$\nu$ 为光子频率），而 $\nu = \frac{c}{\lambda}$，所以 $E = h\frac{c}{\lambda}$，故 A 正确。
}

\item
%% number: 7
%% paperName: 2010年上海高考第 7 题 3 分
%% typeId: 1
%% type: 单选题
%% chapter: 光学
%% point: 光的电磁说
%% method: 无
%% score: 3
%% degree: 850
%% duplicateId: 0
%% body:
电磁波包含了 $\gamma$ 射线、红外线、紫外线、无线电波等，按波长由长到短的排列顺序是\xzanswer{A}
\fourchoices[answer=A]
{无线电波、红外线、紫外线、$\gamma$ 射线}
{红外线、无线电波、$\gamma$ 射线、紫外线}
{$\gamma$ 射线、红外线、紫外线、无线电波}
{紫外线、无线电波、$\gamma$ 射线、红外线}
%% answer: A
%% memo:
\memoanswer{
在题干中的几种电磁波中，无线电波波长最长，其次是红外线，$\gamma$ 射线最短，故 A 正确。
}

\item
%% number: 8
%% paperName: 2010年上海高考第 8 题 2 分
%% typeId: 1
%% type: 单选题
%% chapter: 原子和原子核
%% point: 半衰期
%% method: 无
%% score: 2
%% degree: 650
%% duplicateId: 0
%% body:
某放射性元素经过 $11.4\Ud$ 有 $\frac{7}{8}$ 的原子核发生了衰变，该元素的半衰期为\xzanswer{D}
\fourchoices[answer=D]
{$11.4\Ud$}
{$7.6\Ud$}
{$5.7\Ud$}
{$3.8\Ud$}
%% answer: D
%% memo:
\memoanswer{
这种元素经 $11.4\Ud$，还剩余 $\frac{1}{8}$，即经过了 3 个半衰期，故半衰期为 $\frac{11.4}{3}\Ud=3.8\Ud$，故 D 正确。
}

\item
%% number: 9
%% paperName: 2010年上海高考第 9 题 3 分
%% typeId: 1
%% type: 单选题
%% chapter: 电场
%% point: 电场线
%% method: 无
%% score: 3
%% degree: 650
%% duplicateId: 0
%% body:
三个点电荷电场的电场线分布如图所示，图中 a、b 两点处的场强大小分别为 $E_{a}$、$E_{b}$，电势分别为 $\varphi_{a}$、$\varphi_{b}$，则\xzanswer{C}
\onepicture[w=4.5cm]{\includesvg{figs/09}}
\fourchoices[answer=C]
{$E_{a} > E_{b}$，$\varphi_{a} > \varphi_{b}$}
{$E_{a} < E_{b}$，$\varphi_{a} < \varphi_{b}$}
{$E_{a} > E_{b}$，$\varphi_{a} < \varphi_{b}$}
{$E_{a} < E_{b}$，$\varphi_{a} > \varphi_{b}$}
%% answer: C
%% memo:
\memoanswer{
由电场线的疏密程度可以判断，a 点场强大于 b 点场强；沿电场线电势逐渐降低，作出 a 点等势面会发现 a 点的电势低于 b 点，综上所述 $E_{a}>E_{b},\varphi_{a}<\varphi_{b}$，故 C 正确。
}

\item
%% number: 10
%% paperName: 2010年上海高考第 10 题 3 分
%% typeId: 1
%% type: 单选题
%% chapter: 气体性质
%% point: 液柱移动定性分析
%% method: 无
%% score: 3
%% degree: 650
%% duplicateId: 0
%% body:
如图，玻璃管内封闭了一段气体，气柱长度为 $l$，管内外水银面高度差为 $h$。若温度保持不变，把玻璃管稍向上提起一段距离，则\xzanswer{A}
\onepicture[h=4.0cm]{\includesvg{figs/10}}
\fourchoices[answer=A]
{$h$、$l$ 均变大}
{$h$、$l$ 均变小}
{$h$ 变大 $l$ 变小}
{$h$ 变小 $l$ 变大}
%% answer: A
%% memo:
\memoanswer{
假设 $h$ 不变，则随着玻璃管的上提，密封气体的体积增大，由玻意耳定律知，气体的压强减小，因此玻璃管中液面会升高；液面升高后，由于气体压强减小了，所以气体体积会增大，即 $h$ 和 $l$ 都增大，故 A 正确。
}

\item
%% number: 11
%% paperName: 2010年上海高考第 11 题 3 分
%% typeId: 1
%% type: 单选题
%% chapter: 运动和力
%% point: 牛顿定律的应用
%% method: 无
%% score: 3
%% degree: 650
%% duplicateId: 0
%% body:
将一个物体以某一速度从地面竖直向上抛出，设物体在运动过程中所受空气阻力大小不变，则物体\xzanswer{A}
\fourchoices[answer=A]
{刚抛出时的速度最大}
{在最高点的加速度为零}
{上升时间大于下落时间}
{上升时的加速度等于下落时的加速度}
%% answer: A
%% memo:
\memoanswer{
上升过程中 $mg + f = ma_{1}$，下降过程中 $mg - f = ma_{2}$，所以上升过程中的加速度 $a_{1}$ 大于下降过程中的加速度 $a_{2}$，故 D 错误；再由动能定理判断，上升时的初速度 $v_{0}$ 大于落回原处的末速度 $v$，故 A 正确；由 $h = \frac{v}{2}t$ 判断，上升的时间小于下降过程中的时间，故 C 错误；在最高点物体只受重力作用，所以加速度不为零，故 B 错误。
}

\item
%% number: 12
%% paperName: 2010年上海高考第 12 题 3 分
%% typeId: 1
%% type: 单选题
%% chapter: 直线运动
%% point: 运动的合成
%% method: 无
%% score: 3
%% degree: 650
%% duplicateId: 0
%% body:
降落伞在匀速下落过程中遇到水平方向吹来的风，若风速越大，则降落伞\xzanswer{D}
\fourchoices[answer=D]
{下落的时间越短}
{下落的时间越长}
{落地时速度越小}
{落地时速度越大}
%% answer: D
%% memo:
\memoanswer{
由运动的独立性知，水平方向吹来的风不影响竖直方向的运动，所以下落的时间不变，落地时的竖直分速度也不变；落地时的速度等于竖直分速度和水平分速度的矢量和 $v=\sqrt{v_{y}^{2}+v_{x}^{2}}$，当 $v_{x}$ 越大时，则落地时的速度越大，故 D 正确。
}

\item
%% number: 13
%% paperName: 2010年上海高考第 13 题 3 分
%% typeId: 1
%% type: 单选题
%% chapter: 磁场
%% point: 安培力
%% method: 等效替代
%% score: 3
%% degree: 650
%% duplicateId: 0
%% body:
如图，长为 $2l$ 的直导线折成边长相等，夹角为 $60^{\circ}$ 的 V 形，并置于与其所在平面相垂直的匀强磁场中，磁感应强度为 $B$。当在该导线中通以电流强度为 $I$ 的电流时，该 V 形通电导线受到的安培力大小为\xzanswer{C}
\onepicture[w=4.0cm]{\includesvg{figs/13}}
\fourchoices[answer=C]
{0}
{$0.5BIl$}
{$BIl$}
{$2BIl$}
%% answer: C
%% memo:
\memoanswer{
安培力公式 $F = BIL$，其中 $L$ 为垂直放于磁场中导线的有效长度，图中的有效长度为虚线所示，其长度为 $l$，所以通电导线受到的安培力 $F = BIl$，故 C 正确。

\onepicture[h=3.5cm, align=right]{figs/13_an.jpg}
}

\item
%% number: 14
%% paperName: 2010年上海高考第 14 题 3 分
%% typeId: 1
%% type: 单选题
%% chapter: 气体性质
%% point: 分子力
%% method: 无
%% score: 3
%% degree: 650
%% duplicateId: 0
%% body:
分子间的相互作用力由引力与斥力共同产生，并随着分子间距离的变化而变化。则\xzanswer{B}
\fourchoices[answer=B]
{分子间引力随分子间距的增大而增大}
{分子间斥力随分子间距的减小而增大}
{分子间相互作用力随分子间距的增大而增大}
{分子间相互作用力随分子间距的减小而增大}
%% answer: B
%% memo:
\memoanswer{
分子力随分子间距离变化的关系如图所示。从图象中可以看出，分子间斥力和引力都随分子间距离的增大而减小，随分子间距离的减小而增大，但斥力变化得更快一些，故 A 错误、B 正确；实际的分子力是引力和斥力的合力（如图中实线所示），可以看出，在 $r > r_{0}$ 时，随着分子间距离的减小，分子力先增大后减小，当 $r = r_{0}$ 时分子力为零，故 CD 均错误。

\onepicture[h=3.6cm, align=right]{figs/14_an.jpg}
}

\item
%% number: 15
%% paperName: 2010年上海高考第 15 题 3 分
%% typeId: 1
%% type: 单选题
%% chapter: 曲线运动
%% point: 万有引力定律
%% method: 无
%% score: 3
%% degree: 650
%% duplicateId: 0
%% body:
月球绕地球做匀速圆周运动的向心加速度大小为 $a$，设月球表面的重力加速度大小为 $g_{1}$，在月球绕地球运行的轨道处由地球引力产生的重力加速度为 $g_{2}$。则\xzanswer{B}
\fourchoices[answer=B]
{$g_{1} = a$}
{$g_{2} = a$}
{$g_{1} + g_{2} = a$}
{$g_{2} - g_{1} = a$}
%% answer: B
%% memo:
\memoanswer{
设月球绕地球做匀速圆周运动的半径为 $R$，地球的质量为 $M$，月球的质量为 $m$，由万有引力定律得 $G\frac{Mm}{R^{2}}=ma=mg_{2}$；设月球的半径为 $r$，月球表面有一重物质量为 $\Delta m$，由万有引力得 $G\frac{m\Delta m}{r^{2}}=\Delta mg_{1}$，故 B 正确。
}

\item
%% number: 16
%% paperName: 2010年上海高考第 16 题 3 分
%% typeId: 1
%% type: 单选题
%% chapter: 机械波
%% point: 波的图象
%% method: 无
%% score: 3
%% degree: 650
%% duplicateId: 0
%% body:
如右图，一列简谐横波沿 $x$ 轴正方向传播，实线和虚线分别表示 $t_{1} = 0$ 和 $t_{2} = 0.5\Us$（$T > 0.5\Us$）时的波形，能正确反映 $t_{3} = 7.5\Us$ 时波形的是图\xzanswer{D}
\onepicture[w=5.0cm]{\includesvg{figs/16a}}
\fourchoices[answer=D, ispicture=true, h=2.0cm]
{figs/16b.png}
{figs/16c.png}
{figs/16d.png}
{figs/16e.png}
%% answer: D
%% memo:
\memoanswer{
由 $\frac{1}{4}T=0.5\Us$ 得 $T=2\Us$，则 $n=\frac{t}{T}=\frac{7.5}{2}=3\frac{3}{4}$，经 $3\frac{3}{4}T$，$x=0$ 的质点此时运动到波峰处，故 D 正确。
}

\item
%% number: 17
%% paperName: 2010年上海高考第 17 题 4 分
%% typeId: 2
%% type: 多选题
%% chapter: 气体性质
%% point: 气体图象
%% method: 无
%% score: 4
%% degree: 650
%% duplicateId: 0
%% body:
一定量的理想气体的状态经历了如图所示的 ab、bc、cd、da 四个过程。其中 bc 的延长线通过原点，cd 垂直于 ab 且与水平轴平行，da 和 bc 平行。则气体体积在\xzanswer{AB}
\onepicture[w=4.5cm]{\includesvg{figs/17}}
\fourchoices[answer=AB]
{ab 过程中不断增加}
{bc 过程中保持不变}
{cd 过程中不断增加}
{da 过程中保持不变}
%% answer: AB
%% memo:
\memoanswer{
ab 过程为等温过程，压强减小则体积增加，故 A 正确；

bc 为等容过程，体积不变，故 B 正确；

cd 为等压过程，温度降低则体积减小，故 C 错误；

如图所示，连接 aO 交 cd 于 e，则 ae 是等容线，即 $V_{a}=V_{e}$，因为 $V_{d}<V_{e}$，所以 $V_{d}<V_{a}$，即 da 过程中体积不是保持不变，选项 D 错误。

\onepicture[h=4.0cm, align=right]{figs/17_an.png}
}

\item
%% number: 18
%% paperName: 2010年上海高考第 18 题 4 分
%% typeId: 2
%% type: 多选题
%% chapter: 直线运动
%% point: 运动图象
%% method: 无
%% score: 4
%% degree: 650
%% duplicateId: 0
%% body:
如图为质量相等的两个质点 A、B 在同一直线上运动的 $v-t$ 图像。由图可知\xzanswer{BCD}
\onepicture[w=4.5cm]{\includesvg{figs/18}}
\fourchoices[answer=BCD]
{在 $t$ 时刻两个质点在同一位置}
{在 $t$ 时刻两个质点速度相等}
{在 $0 \sim t$ 时间内质点 B 比质点 A 位移大}
{在 $0 \sim t$ 时间内合外力对两个质点做功相等}
%% answer: BCD
%% memo:
\memoanswer{
在 $t$ 时刻，两质点的速度相等，但不在同一位置，故 B 正确；$v-t$ 图线与 $t$ 轴所围的面积在数值上等于质点发生的位移，故 C 正确；由于不知两质点的初始位置关系，可知 A 错误；由动能定理知 D 正确。
}

\item
%% number: 19
%% paperName: 2010年上海高考第 19 题 4 分
%% typeId: 2
%% type: 多选题
%% chapter: 电磁感应
%% point: 电磁感应中图象
%% method: 无
%% score: 4
%% degree: 450
%% duplicateId: 0
%% body:
如右图，一有界区域内，存在着磁感应强度大小均为 $B$，方向分别垂直于光滑水平桌面向下和向上的匀强磁场，磁场宽度均为 $L$，边长为 $L$ 的正方形线框 abcd 的 bc 边紧靠磁场边缘置于桌面上，使线框从静止开始沿 $x$ 轴正方向匀加速通过磁场区域，若以逆时针方向为电流的正方向，能反映线框中感应电流变化规律的是图\xzanswer{AC}
\onepicture[w=4.5cm]{\includesvg{figs/19a}}
\fourchoices[answer=AC, ispicture=true, h=2.0cm]
{figs/19b.png}
{figs/19c.png}
{figs/19d.png}
{figs/19e.png}
%% answer: AC
%% memo:
\memoanswer{
感应电流大小为 $I=\frac{BLv}{R}=\frac{BLat}{R}=\frac{BL\sqrt{2as}}{R}$，$I$ 与时间 $t$ 成线性关系，当 bc 边进入第二个磁场后，bc 和 ad 边均切割磁感线，回路中的感应电动势是进入前的 2 倍，但感应电流的方向相反；bc 边开始出第二个磁场时，只有一条边切割磁感线，电动势为出来之前瞬间的 $\frac{1}{2}$ 倍，故 A 正确、B 错误；由 $I$ 与 $x$ 的关系，同理可知 C 正确、D 错误。
}

\item
%% number: 20
%% paperName: 2010年上海高考第 20 题 4 分
%% typeId: 2
%% type: 多选题
%% chapter: 机械波
%% point: 波的多解性
%% method: 无
%% score: 4
%% degree: 450
%% duplicateId: 0
%% body:
如图，一列沿 $x$ 轴正方向传播的简谐横波，振幅为 $2\Ucm$，波速为 $2\Ums$。在波的传播方向上两质点 a、b 的平衡位置相距 $0.4\Um$（小于一个波长），当质点 a 在波峰位置时，质点 b 在 $x$ 轴下方与 $x$ 轴相距 $1\Ucm$ 的位置。则\xzanswer{ACD}
\onepicture[w=6.0cm]{\includesvg{figs/20}}
\fourchoices[answer=ACD]
{此波的周期可能为 $0.6\Us$}
{此波的周期可能为 $1.2\Us$}
{从此时刻起经过 $0.5\Us$，b 点可能在波谷位置}
{从此时刻起经过 $0.5\Us$，b 点可能在波峰位置}
%% answer: ACD
%% memo:
\memoanswer{
由题意知，a、b 所处的位置如图所示，b 点对应的位置有两种情况，第一种情况如 $b_{1}$，则 $\frac{1}{4}\lambda+\frac{1}{12}\lambda=0.4\Um$，解得 $\lambda=1.2\Um$，则周期为 $T=\frac{\lambda}{v}=\frac{1.2}{2}\Us=0.6\Us$；从此时刻起经过 $0.5\Us$，即 $\frac{5}{6}T$，波沿 $x$ 轴正方向传播 $\frac{5}{6}\lambda = 1.0\Um$，波峰到 $x = 1.2\Um$ 处，b 在波谷，故 AC 正确；第二种情况如 $b_{2}$，$\frac{3}{4}\lambda - \frac{1}{12}\lambda = 0.4\Um$，解得 $\lambda = 0.6\Um$，则周期为 $T = \frac{\lambda}{v} = \frac{0.6}{2}\Us = 0.3\Us$，从此时刻起经过 $0.5\Us$，即 $\frac{5}{3}T$，波沿 $x$ 轴正方向传播 $\frac{5}{3}\lambda = 1.0\Um$，波峰到 $x = 1.0\Um$ 处，$x = 0.4\Um$ 的 b 在波峰，故 B 错误、D 正确。

\onepicture[h=3.2cm, align=right]{figs/20_an.jpg}
}

\item
%% number: 21
%% paperName: 2010年上海高考第 21 题 4 分
%% typeId: 3
%% type: 填空题
%% chapter: 电磁感应
%% point: 楞次定律推广
%% method: 无
%% score: 4
%% degree: 650
%% duplicateId: 0
%% body:
如图，金属环 A 用轻线悬挂，与长直螺线管共轴，并位于其左侧。若变阻器滑片 P 向左移动，则金属环 A 将向\tkanswer[1.4cm]{左}（填“左”或“右”）运动，并有\tkanswer[1.6cm]{收缩}（填“收缩”或“扩张”）趋势。
\onepicture[align=right, w=4.0cm]{\includesvg{figs/21}}
%% answer: 左，收缩
%% memo:
\memoanswer{
当变阻器滑片 P 向左移动时，通过螺线管中的电流增大，它所产生的磁场增强，通过环 A 的磁通量增加，由楞次定律知，环 A 中产生的感应电流将阻碍磁通量的增大，故环 A 向左运动，并有收缩的趋势。
}

\item
%% number: 22
%% paperName: 2010年上海高考第 22 题 4 分
%% typeId: 3
%% type: 填空题
%% chapter: 气体性质
%% point: 盖吕萨克定律
%% method: 无
%% score: 4
%% degree: 850
%% duplicateId: 0
%% body:
如图，上端开口的圆柱形气缸竖直放置，截面积为 $5\times10^{-3}\Umq$，一定质量的气体被质量为 $2.0\Ukg$ 的光滑活塞封闭在气缸内，其压强为\tkanswer[2.6cm]{$1.05\times10^{5}$}$\UPa$（大气压强取 $1.01\times10^{5}\UPa$，$g$ 取 $10\Umsq$）。若从初温 $27\UdC$ 开始加热气体，使活塞离气缸底部的高度由 $0.50\Um$ 缓慢地变为 $0.51\Um$。则此时气体的温度为\tkanswer[1.6cm]{$33$}$\UdC$。
\onepicture[align=right, h=3.2cm]{\includesvg{figs/22}}
%% answer: $1.05\times10^{5}$，33
%% memo:
\memoanswer{
气体的压强为 $p = p_{0} + \frac{mg}{S} = 1.01 \times 10^{5}\UPa + \frac{2.0 \times 10}{5 \times 10^{-3}}\UPa = 1.05 \times 10^{5}\UPa$；加热过程为等压过程，由 $\frac{h_{2}S}{h_{1}S} = \frac{T_{2}}{T_{1}}$，解得 $T_{2} = \frac{h_{2}}{h_{1}}T_{1} = \frac{51}{50} \times 300\UK = 306\UK$，则 $t_{2} = (306 - 273)\UdC = 33\UdC$。
}

\item
%% number: 23
%% paperName: 2010年上海高考第 23 题 4 分
%% typeId: 3
%% type: 填空题
%% chapter: 电路
%% point: 逻辑电路
%% method: 无
%% score: 4
%% degree: 650
%% duplicateId: 0
%% body:
电动机的自动控制电路如图所示，其中 $R_{H}$ 为热敏电阻，$R_{L}$ 为光敏电阻，当温度升高时，$R_{H}$ 的阻值远小于 $R_{1}$；当光照射 $R_{L}$ 时，其阻值远小于 $R_{2}$，为使电动机在温度升高或受到光照时能自动启动，电路中虚线框内应选\tkanswer[1.6cm]{或}门逻辑电路；若要提高光照时电动机启动的灵敏度，可以\tkanswer[1.6cm]{增大}$R_{2}$ 的阻值（填“增大”或“减小”）。
\onepicture[align=right, w=5.8cm]{\includesvg{figs/23}}
%% answer: 或，增大
%% memo:
\memoanswer{
无论 $R_{H}$ 还是 $R_{L}$ 阻值较小时，对门电路输入的都是高电势，因此框内为或门电路；要提高光照时电动机启动的灵敏度，就是提高 $R_{L}$ 输出端的电势，因此可以增大 $R_{2}$ 的阻值。
}

\item
%% number: 24
%% paperName: 2010年上海高考第 24 题 4 分
%% typeId: 3
%% type: 填空题
%% chapter: 曲线运动
%% point: 天体运动
%% method: 无
%% score: 4
%% degree: 650
%% duplicateId: 0
%% body:
如图，三个质点 a、b、c 质量分别为 $m_{1}$、$m_{2}$、$M$（$M \gg m_{1}$，$M \gg m_{2}$），在 c 的万有引力作用下，a、b 在同一平面内绕 c 沿逆时针方向做匀速圆周运动，轨道半径之比为 $r_{a}:r_{b}=1:4$，则它们的周期之比 $T_{a}:T_{b}=$\tkanswer[2.0cm]{$1:8$}，从图示位置开始，在 b 转动一周的过程中，a、b、c 共线有\tkanswer[1.2cm]{$14$}次。
\onepicture[align=right, w=4.8cm]{\includesvg{figs/24}}
%% answer: $1:8$，14
%% memo:
\memoanswer{
c 对 a、b 作用的万有引力提供它们绕 c 做圆周运动的向心力，由牛顿第二定律得 $G\frac{Mm}{r^{2}}=m\frac{4\pi^{2}}{T^{2}}r$，得 $\frac{r^{3}}{T^{2}}=$ 常数，所以有 $\frac{r_{a}^{3}}{r_{b}^{3}}=\frac{T_{a}^{2}}{T_{b}^{2}}$，代入数据得 $\frac{T_{a}}{T_{b}}=\frac{1}{8}$；设每隔时间 $t$，a、b 共线一次，则 $(\omega_{a}-\omega_{b})t=\pi$，所以 $t=\frac{\pi}{(\omega_{a}-\omega_{b})}$，所以 b 运动一周的过程中，a、b、c 共线的次数为 $n=\frac{T_{b}}{t}=\frac{T_{b}(\omega_{a}-\omega_{b})}{\pi}=T_{b}\left(\frac{2}{T_{a}}-\frac{2}{T_{b}}\right)=\frac{2T_{b}}{T_{a}}-2=14$。
}

\item
%% number: 25
%% paperName: 2010年上海高考第 25 题 4 分
%% typeId: 3
%% type: 填空题
%% chapter: 功和能
%% point: 动能定理
%% method: 无
%% score: 4
%% degree: 650
%% duplicateId: 0
%% body:
如图，固定于竖直面内的粗糙斜杆，与水平方向夹角为 $30^{\circ}$，质量为 $m$ 的小球套在杆上，在大小不变的拉力作用下，小球沿杆由底端匀速运动到顶端。为使拉力做功最小，拉力 $F$ 与杆的夹角 $\alpha =$\tkanswer[1.8cm]{$60^{\circ}$}，拉力大小 $F =$\tkanswer[1.6cm]{$mg$}。
\onepicture[align=right, w=4.8cm]{\includesvg{figs/25}}
%% answer: $60^{\circ}$，$mg$
%% memo:
\memoanswer{
在球匀速沿杆向上运动的过程中，由动能定理得 $W_{F}-mgh-W_{f}=0$，当摩擦力为零时，摩擦力做的功 $W_{f}=0$，此时拉力做功最小，而摩擦力为 0 时，杆对球的弹力一定也为 0，故球只受重力和拉力作用，两个力为平衡力，拉力的方向竖直向上，与杆的夹角为 $60^{\circ}$，大小为 $mg$。
}

\item
%% number: 26
%% paperName: 2010年上海高考第 26 题 6 分
%% typeId: 4
%% type: 实验题
%% chapter: 电场
%% point: 描绘等势面
%% method: 无
%% score: 6
%% degree: 850
%% duplicateId: 0
%% body:
在 DIS 描绘电场等势线的实验中
\begin{enumerate}
	\item
	电源通过正负电极 a、b 在导电物质上产生的稳定电流分布模拟了由二个\tkanswer[3.4cm]{等量异种点电荷}产生的\tkanswer[2.0cm]{静电场}，用\tkanswer[2.2cm]{电压传感器}探测等势点。
	\item
	（单选题）在安装实验装置时，正确的做法是（\quad）

	（A）在一块平整木板上依次铺放复写纸、白纸、导电纸

	（B）导电纸有导电物质的一面应该向上

	（C）连接电源正负极的电极 a、b 必须与导电物质保持绝缘

	（D）连接电极 a、b 的电源电压为交流 $4 \sim 6\UV$
\end{enumerate}
%% answer: （1）等量异种点电荷，静电场，电压传感器（2）B
\jdanswer{
\begin{enumerate}
	\item
	等量异种点电荷，静电场，电压传感器

	\item
	B
\end{enumerate}
}
%% memo:
\memoanswer{
（1）点电荷，电场，电压传感器。

（2）平整木板从下往上依次铺放白纸、复写纸、导电纸，有导电物质的一面应该朝上，故 B 正确、A 错误；连接电源正负极的电极必须与导电物质保持良好接触，故 C 错误；电源应为直流学生电源，故 D 错误。
}

\item
%% number: 27
%% paperName: 2010年上海高考第 27 题 6 分
%% typeId: 4
%% type: 实验题
%% chapter: 曲线运动
%% point: 万有引力定律
%% method: 无
%% score: 6
%% degree: 650
%% duplicateId: 0
%% body:
卡文迪许利用如图所示的扭秤实验装置测量了引力常量 $G$。
\begin{enumerate}
	\item
	（多选题）为了测量石英丝极微小的扭转角，该实验装置中采取使“微小量放大”的主要措施是（\quad）

	（A）减小石英丝的直径

	（B）增大 T 型架横梁的长度

	（C）利用平面镜对光线的反射

	（D）增大刻度尺与平面镜的距离
	\item
	已知 T 型架水平横梁长度为 $l$，质量分别为 $m$ 和 $m'$ 的球，位于同一水平面内，当横梁处于力矩平衡状态时，测得 $m$、$m'$ 连线长度为 $r$，且与水平横梁垂直，同时测得石英丝的扭转角度为 $\theta$，由此得到扭转力矩 $k\theta$（$k$ 为扭转系数且已知），则引力常量的表达式 $G =$\tkanswer[2.4cm]{$\frac{k\theta r^{2}}{mm'l}$}。
\end{enumerate}
\onepicture[align=right, w=7.4cm]{\includesvg{figs/27}}
%% answer: （1）CD（2）$\frac{k\theta r^{2}}{mm'l}$
\jdanswer{
\begin{enumerate}
	\item
	CD

	\item
	$\frac{k\theta r^{2}}{mm'l}$
\end{enumerate}
}
%% memo:
\memoanswer{
（1）利用平面镜对光线反射和增大刻度尺与平面镜的距离可清晰地看到光点在尺子上移动的距离，故 CD 正确。

（2）由力矩平衡得 $2G\frac{mm'}{r^{2}}\cdot\frac{l}{2}=k\theta$，解得 $G=\frac{k\theta r^{2}}{mm'l}$。
}

\item
%% number: 28
%% paperName: 2010年上海高考第 28 题 6 分
%% typeId: 4
%% type: 实验题
%% chapter: 气体性质
%% point: 验证玻意耳定律
%% method: 整体与隔离
%% score: 6
%% degree: 650
%% duplicateId: 0
%% body:
用 DIS 研究一定质量气体在温度不变时，压强与体积关系的实验装置如图~\ref{2010上海28a} 所示，实验步骤如下：
\begin{steps}
	\item
	把注射器活塞移至注射器中间位置，将注射器与压强传感器、数据采集器、计算机逐一连接；
	\item
	移动活塞，记录注射器的刻度值 $V$，同时记录对应的由计算机显示的气体压强值 $p$；
	\item
	用 $V-1/p$ 图像处理实验数据，得到如图~\ref{2010上海28b} 所示图线。
\end{steps}
\begin{enumerate}
	\item
	为了保持封闭气体的质量不变，实验中采取的主要措施是\tkanswer[3.4cm]{在注射器活塞上涂润滑油}。
	\item
	为了保持封闭气体的温度不变，实验中采取的主要措施是\tkanswer[3.0cm]{移动活塞要缓慢}和\tkanswer[3.6cm]{不能用手握住注射器封闭气体部分}。
	\item
	如果实验操作规范正确，但如图所示的 $V-1/p$ 图线不过原点，则 $V_{0}$ 代表\tkanswer[3.6cm]{注射器与压强传感器连接部位的气体体积}。
\end{enumerate}
\twopicture[labelA=2010上海28a, labelB=2010上海28b, numA=1, numB=2, showlabel=false, align=right, wA=5.0cm, wB=4.4cm]
{\includesvg{figs/28a}}
{\includesvg{figs/28b}}
%% answer: （1）在注射器活塞上涂润滑油（2）移动活塞要缓慢；不能用手握住注射器封闭气体部分（3）注射器与压强传感器连接部位的气体体积
\jdanswer{
\begin{enumerate}
	\item
	在注射器活塞上涂润滑油

	\item
	移动活塞要缓慢；不能用手握住注射器封闭气体部分

	\item
	注射器与压强传感器连接部位的气体体积
\end{enumerate}
}
%% memo:
\memoanswer{
（1）在注射器活塞上涂润滑油，防止气体泄漏。

（2）慢慢抽动活塞，目的是使实验过程中的每一微小过程等效于准静态过程。另外，为了防止注射器中气体温度升高，不能用手握住注射器。

（3）体积读数值比实际值小，说明在传感器与注射器之间还有少量的气体，$V_{0}$ 代表的就是该部分气体的体积。
}

\item
%% number: 29
%% paperName: 2010年上海高考第 29 题 6 分
%% typeId: 4
%% type: 实验题
%% chapter: 电路
%% point: 测电动势与内阻
%% method: 无
%% score: 6
%% degree: 850
%% duplicateId: 0
%% body:
某同学利用 DIS、定值电阻 $R_{0}$、电阻箱 $R_{1}$ 等实验器材测量电池 a 的电动势和内阻，实验装置如图~\ref{2010上海29a} 所示。实验时多次改变电阻箱的阻值，记录外电路的总电阻阻值 $R$，用电压传感器测得端电压 $U$，并在计算机上显示出如图~\ref{2010上海29b} 所示的 $\frac{1}{U}-\frac{1}{R}$ 关系图线 a。重复上述实验方法测量电池 b 的电动势和内阻，得到图~\ref{2010上海29b} 中的图线 b。
\begin{enumerate}
	\item
	由图线 a 可知电池 a 的电动势 $E_{a} =$\tkanswer[2.0cm]{$2$}$\UV$，内阻 $r_{a} =$\tkanswer[2.0cm]{$0.5$}$\UO$。
	\item
	若用同一个电阻 $R$ 先后与电池 a 及电池 b 连接，则两电池的输出功率 $P_{a}$\tkanswer[1.6cm]{小于}$P_{b}$（填“大于”、“等于”或“小于”），两电池的效率 $\eta_{a}$\tkanswer[1.6cm]{大于}$\eta_{b}$（填“大于”、“等于”或“小于”）。
\end{enumerate}
\twopicture[labelA=2010上海29a, labelB=2010上海29b, numA=1, numB=2, showlabel=false, align=right, wA=5.8cm, wB=7.2cm]
{\includesvg{figs/29a}}
{\includesvg{figs/29b}}
%% answer: （1）$2\UV$，$0.5\UO$（2）小于，大于
\jdanswer{
\begin{enumerate}
	\item
	$2\UV$，$0.5\UO$

	\item
	小于，大于
\end{enumerate}
}
%% memo:
\memoanswer{
（1）由闭合电路的欧姆定律得 $E=U+\frac{U}{R_{0}+R_{1}}r$，变形得 $\frac{1}{U}=\frac{1}{E}+\frac{r}{E}\cdot\frac{1}{R_{0}+R_{1}}$，对 a 电源，截距 $\frac{1}{E}=0.5$，得到 $E=2.0\UV$；求 $r$ 的方法一：斜率 $\frac{r}{E}=\frac{0.5}{2}$，得 $r=0.5\UO$。求 $r$ 的方法二：由横轴的意义求内阻 $r$，当 $\frac{1}{U}=0$ 时，即 $\frac{1}{E}+\frac{r}{E}\cdot\frac{1}{R_{0}+R_{1}}=0$，进一步得到 $\frac{1}{r}=-\frac{1}{R_{0}+R_{1}}$，所以 $r=\frac{1}{2.0}\UO=0.5\UO$。

（2）从上述分析可以看出，$E$ 等于纵轴截距的倒数，$r$ 等于横轴截距的倒数，斜率 $\frac{r}{E}=\frac{1}{I_{\text{短}}}$，所以有 $E_{b}>E_{a}, r_{b}>r_{a}$，电池的输出功率 $P=\left(\frac{E}{R+r}\right)^{2}R$ 得 $P_{a}<P_{b}$，又由于 $\eta=\frac{UI}{EI}=\frac{R}{R+r}$，则 $\eta_{a}>\eta_{b}$。
}

\item
%% number: 30
%% paperName: 2010年上海高考第 30 题 10 分
%% typeId: 6
%% type: 计算题
%% chapter: 曲线运动
%% point: 平抛运动
%% method: 无
%% score: 10
%% degree: 650
%% duplicateId: 0
%% body:
如图，ABC 和 ABD 为两个光滑固定轨道，A、B、E 在同一水平面上，C、D、E 在同一竖直线上，D 点距水平面的高度为 $h$，C 点的高度为 $2h$，一滑块从 A 点以初速度 $v_{0}$ 分别沿两轨道滑行到 C 或 D 处后水平抛出。
\begin{enumerate}
	\item
	求滑块落到水平面时，落点与 E 点间的距离 $s_{C}$ 和 $s_{D}$；
	\item
	为实现 $s_{C} < s_{D}$，$v_{0}$ 应满足什么条件？
\end{enumerate}
\onepicture[align=right, w=9.0cm]{\includesvg{figs/30}}
%% answer: （1）$s_{C}=\sqrt{v_{0}^{2}-4gh}\sqrt{\frac{4h}{g}}$，$s_{D}=\sqrt{v_{0}^{2}-2gh}\sqrt{\frac{2h}{g}}$（2）$\sqrt{4gh}<v_{0}<\sqrt{6gh}$
\jdanswer{
\begin{enumerate}
	\item
	$s_{C}=\sqrt{v_{0}^{2}-4gh}\sqrt{\frac{4h}{g}}$，$s_{D}=\sqrt{v_{0}^{2}-2gh}\sqrt{\frac{2h}{g}}$

	\item
	$\sqrt{4gh}<v_{0}<\sqrt{6gh}$
\end{enumerate}
}
%% memo:
\memoanswer{
（1）设抛出点高度为 $y$，根据机械能守恒，有

$\frac{1}{2}mv_{0}^{2} = \frac{1}{2}mv^{2} + mgy$，

平抛初速度

$v = \sqrt{v_{0}^{2} - 2gy}$，

落地时间 $t$ 满足

$y = \frac{1}{2}gt^{2}$，解得 $t = \sqrt{\frac{2y}{g}}$。

落地点离抛出点的水平距离

$s = vt = \sqrt{v_{0}^{2} - 2gy}\sqrt{\frac{2y}{g}}$，

分别以 $y = 2h$ 和 $y = h$ 代入得

$s_{C} = \sqrt{v_{0}^{2} - 4gh}\sqrt{\frac{4h}{g}}$，$s_{D} = \sqrt{v_{0}^{2} - 2gh}\sqrt{\frac{2h}{g}}$。

（2）由题意 $s_{C} < s_{D}$，即

$\sqrt{v_{0}^{2} - 4gh}\sqrt{\frac{4h}{g}} < \sqrt{v_{0}^{2} - 2gh}\sqrt{\frac{2h}{g}}$，

得 $v_{0}^{2} < 6gh$。

但滑块从 A 点以初速度 $v_{0}$ 分别沿两轨道滑行到 C 或 D 处后水平抛出，要求滑块能到达 C 点，即

$v_{C}^{2} = v_{0}^{2} - 4gh > 0$，

故 $v_{0}^{2} > 4gh$，

因此初速度应满足 $\sqrt{4gh} < v_{0} < \sqrt{6gh}$。
}

\item
%% number: 31
%% paperName: 2010年上海高考第 31 题 12 分
%% typeId: 6
%% type: 计算题
%% chapter: 运动和力
%% point: 牛顿第二定律
%% method: 整体与隔离
%% score: 12
%% degree: 650
%% duplicateId: 0
%% body:
倾角 $\theta = 37^{\circ}$，质量 $M = 5\Ukg$ 的粗糙斜面位于水平地面上。质量 $m = 2\Ukg$ 的木块置于斜顶端，从静止开始匀加速下滑，经 $t = 2\Us$ 到达底端，运动路程 $L = 4\Um$，在此过程中斜面保持静止（$\sin 37^{\circ} = 0.6$，$\cos 37^{\circ} = 0.8$，$g$ 取 $10\Umsq$）。求：
\begin{enumerate}
	\item
	地面对斜面的摩擦力大小与方向；
	\item
	地面对斜面的支持力大小；
	\item
	通过计算证明木块在此过程中满足动能定理。
\end{enumerate}
\onepicture[align=right, w=5.3cm]{\includesvg{figs/31}}
%% answer: （1）$3.2\UN$，方向水平向左（2）$67.6\UN$（3）见解析
\jdanswer{
\begin{enumerate}
	\item
	$3.2\UN$，方向水平向左

	\item
	$67.6\UN$

	\item
	见解析
\end{enumerate}
}
%% memo:
\memoanswer{
（1）木块做加速运动，有

$L = \frac{1}{2}at^{2}$，解得 $a = \frac{2L}{t^{2}} = 2\Umsq$。

对木块（受力分析图如图所示），由牛顿第二定律

$mg\sin\theta - f_{1} = ma$，

解得 $f_{1} = mg\sin\theta - ma = 8\UN$，

$N_{1} = mg\cos\theta = 16\UN$，

\onepicture[h=3.2cm, align=right]{\includesvg{figs/31_an1}}

斜面受力如图，由共点力平衡，地对斜面的摩擦力

$f_{2} = N_{1}'\sin\theta - f_{1}'\cos\theta = 3.2\UN$，方向水平向左。

\onepicture[h=3.0cm, align=right]{\includesvg{figs/31_an2}}

（2）地面对斜面的支持力为

$N_{2} = Mg + N_{1}\cos\theta + f_{1}\sin\theta = 67.6\UN$。

（3）木块在下滑过程中，沿斜面方向合力及该力做的功为

$F = mg\sin\theta - f_{1} = 4\UN$，

$W = FL = 16\UJ$，

木块末速度及动能增量 $v = at = 4\Ums$，$\Delta E_{k} = \frac{1}{2}mv^{2} = 16\UJ$，

由此可知下滑过程中 $W = \Delta E_{k}$，即合力做功或外力对木块做的总功等于动能的变化（增加），动能定理成立。
}

\item
%% number: 32
%% paperName: 2010年上海高考第 32 题 14 分
%% typeId: 5
%% type: 辨析题
%% chapter: 电磁感应
%% point: 电磁感应中的匀变速
%% method: 无
%% score: 14
%% degree: 450
%% duplicateId: 0
%% body:
如图，宽度为 $L = 0.5\Um$ 的光滑金属框架 MNPQ 固定于水平面内，并处在磁感应强度大小 $B = 0.4\UT$，方向竖直向下的匀强磁场中，框架的电阻非均匀分布。将质量 $m = 0.1\Ukg$，电阻可忽略的金属棒 ab 放置在框架上，并与框架接触良好。以 P 为坐标原点，PQ 方向为 $x$ 轴正方向建立坐标。金属棒从 $x_{0} = 1\Um$ 处以 $v_{0} = 2\Ums$ 的初速度，沿 $x$ 轴负方向做 $a = 2\Umsq$ 的匀减速直线运动，运动中金属棒仅受安培力作用。求：
\begin{enumerate}
	\item
	金属棒 ab 运动 $0.5\Um$，框架产生的焦耳热 $Q$；
	\item
	框架中 aNPb 部分的电阻 $R$ 随金属棒 ab 的位置 $x$ 变化的函数关系；
	\item
	为求金属棒 ab 沿 $x$ 轴负方向运动 $0.4\Us$ 过程中通过 ab 的电量 $q$，某同学解法为：先算出经过 $0.4\Us$ 金属棒的运动距离 $s$，以及 $0.4\Us$ 时回路内的电阻 $R$，然后代入 $q = \frac{\Delta \Phi}{R} = \frac{BLs}{R}$ 求解。指出该同学解法的错误之处，并用正确的方法解出结果。
\end{enumerate}
\onepicture[align=right, w=7.9cm]{\includesvg{figs/32}}
%% answer: （1）$Q=0.1\UJ$（2）$R=0.4\sqrt{x}$（3）错误之处是把 $0.4\Us$ 时回路内的电阻 $R$ 代入 $q=\frac{BLs}{R}$ 进行计算，$q=0.4\UC$
\jdanswer{
\begin{enumerate}
	\item
	$Q=0.1\UJ$

	\item
	$R=0.4\sqrt{x}$

	\item
	错误之处是把 $0.4\Us$ 时回路内的电阻 $R$ 代入 $q=\frac{BLs}{R}$ 进行计算，$q=0.4\UC$
\end{enumerate}
}
%% memo:
\memoanswer{
（1）金属棒仅受安培力作用，其大小

$F = ma = 0.2\UN$。

金属棒运动 $0.5\Um$，框架中产生的焦耳热等于克服安培力做的功，故

$Q = Fs = 0.1\UJ$。

（2）金属棒所受安培力为

$F = BIL$，

又 $I = \frac{E}{R} = \frac{BLv}{R}$，

由牛顿第二定律得

$F = \frac{B^{2}L^{2}v}{R} = ma$，

由于棒做匀减速运动，有

$v = \sqrt{v_{0}^{2} - 2a(x_{0} - x)}$，

联立解得 $R = \frac{B^{2}L^{2}}{ma}\sqrt{v_{0}^{2} - 2a(x_{0} - x)} = 0.4\sqrt{x}$。

（3）错误之处是把 $0.4\Us$ 时回路内的电阻 $R$ 代入 $q = \frac{BLs}{R}$ 进行计算。

正确解法是：由电荷量的定义知

$q = It$，

又 $F = BIL = ma$，

联立解得 $q = \frac{mat}{BL} = 0.4\UC$。
}

\item
%% number: 33
%% paperName: 2010年上海高考第 33 题 14 分
%% typeId: 6
%% type: 计算题
%% chapter: 气体性质
%% point: 气体连接体
%% method: 无
%% score: 14
%% degree: 450
%% duplicateId: 0
%% body:
如图，一质量不计、可上下移动的活塞将圆筒分为上下两室，两室中分别封有理想气体。筒的侧壁为绝缘体，上底 N、下底 M 及活塞 D 均为导体并按图连接，活塞面积 $S = 2\Ucmq$。在电键 K 断开时，两室中气体压强均为 $p_{0} = 240\UPa$，ND 间距 $l_{1} = 1\Uum$，DM 间距 $l_{2} = 3\Uum$。将变阻器的滑片 P 滑到左端 B，闭合电键后，活塞 D 与下底 M 分别带有等量异种电荷，并各自产生匀强电场，在电场力作用下活塞 D 发生移动。稳定后，ND 间距 $l_{1}' = 3\Uum$，DM 间距 $l_{2}' = 1\Uum$，活塞 D 所带电量的绝对值 $q = \varepsilon_{0}SE$（式中 $E$ 为 D 与 M 所带电荷产生的合场强，常量 $\varepsilon_{0} = 8.85\times10^{-12}\UCqNmq$）。求：
\begin{enumerate}
	\item
	两室中气体的压强（设活塞移动前后气体温度保持不变）；
	\item
	活塞受到的电场力大小 $F$；
	\item
	M 所带电荷产生的场强大小 $E_{M}$ 和电源电压 $U$；
	\item
	使滑片 P 缓慢地由 B 向 A 滑动，活塞如何运动，并说明理由。
\end{enumerate}
\onepicture[align=right, w=6.4cm]{\includesvg{figs/33}}
%% answer: （1）$p_{1}=80\UPa$，$p_{2}=720\UPa$（2）$F=0.128\UN$（3）$E_{M}=6\times10^{6}\UVm$，$U=12\UV$（4）见解析
\jdanswer{
\begin{enumerate}
	\item
	$p_{1}=80\UPa$，$p_{2}=720\UPa$

	\item
	$F=0.128\UN$

	\item
	$E_{M}=6\times10^{6}\UVm$，$U=12\UV$

	\item
	当滑片 P 由 B 向 A 滑动时，DM 间场强减小，DN 间场强变大，活塞受到向下的电场力减小，电场力与气体压力间的平衡被破坏，活塞向上移动。
\end{enumerate}
}
%% memo:
\memoanswer{
（1）电键未合上时两室中气体压强为 $p_{0}$，设电键合上后，两室中气体压强分别为 $p_{1}$、$p_{2}$，由玻意耳定律得

$p_{0}l_{1}S = p_{1}l_{1}'S$，$p_{1} = \frac{p_{0}}{3} = 80\UPa$，

$p_{0}l_{2}S = p_{2}l_{2}'S$，$p_{2} = 3p_{0} = 720\UPa$。

（2）活塞受到的气体压强差为

$\Delta p = p_{2} - p_{1} = 640\UPa$，

活塞在气体压力和电场力作用下处于平衡，电场力

$F = \Delta pS = 0.128\UN$。

（3）活塞受到的电场力大小为

$F = qE_{M}$，

其中活塞带电量 $q = \varepsilon_{0}SE$，因为 $E$ 由 D、M 所带等量异号电荷共同产生，根据电场叠加原理，M 产生的场强大小

$E_{M} = \frac{E}{2}$，

则 $F = qE_{M} = 2\varepsilon_{0}SE_{M}^{2}$，

解得 $E_{M} = \sqrt{\frac{F}{2\varepsilon_{0}S}} = \sqrt{\frac{0.128}{2 \times 8.85 \times 10^{-12} \times 2 \times 10^{-4}}}\UVm = 6\times10^{6}\UVm$，

故 $E = 2E_{M} = 1.2\times10^{7}\UVm$，$U = El_{2}' = 1.2\times10^{7} \times 1 \times 10^{-6}\UV = 12\UV$。

（4）当滑片 P 由 B 向 A 滑动时，DM 间场强减小，DN 间场强变大，活塞受到向下的电场力减小，电场力与气体压力间的平衡被破坏，活塞向上移动。
}

\end{enumerate}

\end{document}
